% Transcription — fonds Alexandre Grothendieck, shelfmark 45, batch 4 (pages 61-79).
% Established from the facsimile archives/batches/45/batch-04.pdf, cut from a
% copy of 45.pdf fetched through the project's relay
% (https://grothendieck-relay.onrender.com/source/45.pdf); PDF page k =
% archivists' page 60+k, no cover sheet. Per the skill transcribe-grothendieck.
% The folder is seventy-nine pages in four batches (1-20, 21-40, 41-60,
% 61-79); this is the fourth and last, the others are done by separate passes.
%
% Pass: Opus 5.5 (claude-opus-5-5), 2026-09-26 — first pass, unchecked against the pages by a human.
%
% Pages skipped, without description: 62, 64, 66, 68, 70, 79. They are
% stencilled pages (106, 123, 108, 109, 122, 105) of a text numbered
% « n° 371 » on the exterior algebra (produit intérieur, puissance extérieure
% d'une matrice, « Rabiot qui doit remonter au §3 »), reused as scrap paper:
% each is the verso of one of his leaves (the stencil shows through 61, 63,
% 65, ...). Their underlinings are stencilled; they carry nothing in his hand
% (on 79 a single short ink stroke at the top edge, not a sign). Page 76 is
% transcribed for its top third only; the rest of the leaf is show-through
% of page 75.
%
% THE 1961 ANNOTATED TYPESCRIPT is pages 71-74. It is not a published text
% nor another author's: it is a typed four-page commentary headed
% « Commentaires sur l'article de A. Delzant sur les formes quadratiques »,
% dated in ink top right of page 71 « 22.7.961 » (the second digit written
% over another). Evidence that it is Grothendieck's own: it speaks in the
% first person of « mon séminaire » for the « groupe fondamental » of
% Spec(A) (SGA 1, 1960-61, is his seminar); it addresses Delzant as « tu »
% and prescribes « du point de vue math. ... comme programme minimum »,
% citing « Serre suggère »; and its corrections are in his hand (the Greek
% letters phi, xi, lambda, mu and the p of A/pA inked into blanks the
% typewriter left, « même » for « en effet », « d'anneaux augmentés » for
% « additif », « quadratique(s) », « (pré)schéma », « d'un point s »,
% « S » after « de type fini »), the same blue ink as the date. It is a
% referee's letter on Delzant's article, not the article; no published
% version of it is known to this pass. It is therefore transcribed in full,
% as his unpublished
% text, with his ink interventions marked. The typist's own x-outs
% (« xxx » overstrikes) are omitted silently; an underlined typed Z is set
% \mathbb{Z}, typed underlining of words as \emph.
%
% Pages transcribed: 61, 63, 65, 67, 69 (his notes on the stencil versos),
% 71-74 (the typescript), 75, 76 (top), 77, 78. Pages summarised: none
% wholly; page 61's cancelled upper diagram, page 65's sideways column of
% computations and page 78's struck-through block are given in their legible
% parts, with a \note each.
%
% His pagination: none on 61-69; pages 75-77 carry his circled section
% numbers (1) (75), (2) and (3) (77), set as (1), (2), (3) with a \note; the
% typescript is unpaginated.
%
% Notation, kept from batch 3 where it applies: \underline{\mathrm{Aut}}
% etc. New in this batch: E(A), F(A) the rings of classes of quadratic
% forms (non-degenerate / unimodular) after Kneser, epsilon, d, c rank,
% discriminant, Clifford class, the filtration F_1 ⊃ F_2 ⊃ F_3; Q(A),
% D(A) = A*/A*^2 = H^1(A, Z/2Z), B = H^2(A, Z/2Z), phi_a, e_a, W^i
% Stiefel-Whitney classes in the typescript; E^{n,m}(S) on page 69; mu_2,
% alpha_2, G_m, Pin, Spin; Symm^2, Sigma^2, Alt^2, E^(2) on page 67.
%
% What the batch is about: the run « Formes quadratiques, bilinéaires
% symétriques, etc » named by batch 3's last cover. Page 61: a diagram of
% the categories Bil, Quad, Alt, Mod (typ., unimod., spéciales). Page 63:
% O(2m), O(2m+1), SO and their centres, Pin(m) and Spin(m), and the
% orthogonal group in characteristic 2 in odd rank as mu_2 ⋉ alpha_2^n.
% Page 65: bilinear, alternating, quadratic forms over a ring, H^1(S, mu_2),
% Q(x∧y). Page 67: the exact sequences linking Sigma^2, Symm^2, Alt^2 and
% the Frobenius twist E^(2); tensor products of forms. Page 69: monoids
% E^{n,m}(S) of regular quadratic forms of type (n,m). Pages 71-74: the
% typescript — structure of Q(A) over a field of char ≠ 2 (generators
% phi_a, relations from Witt, grad Q(A) and Symm(D)), Stiefel-Whitney
% classes W^i : Q(A) -> H^i(A, Z/2Z), the local case, comparison with the
% topological classes by GAGA, and char 2 (mu_2 against Z/2Z, orthogonal
% group schemes over Z). Pages 75-77: notes « Exposés de Kneser à l'IHES »:
% E(A), F(A), rank, discriminant, Clifford class, the filtration and its
% quotients, examples (alg. closed, finite, p-adic, R, complete valued
% fields), number fields and Hasse. Page 78: local-global exact sequences
% for E and F over a Dedekind ring, H^1(K, Z/2Z), idèles.

\documentclass[11pt,a4paper]{article}
% Le préambule commun à toutes les transcriptions du fonds Grothendieck.
%
% Il tient en une idée : **l'incertitude se compose, elle ne se résout pas.**
% Une transcription qui lisse un mot douteux a détruit la seule chose pour
% laquelle on la faisait — un lecteur doit pouvoir savoir, à chaque endroit,
% ce qui était sur le feuillet et ce qui a été ajouté. Les six macros qui
% suivent sont donc l'appareil critique, et non de la décoration.
%
% Les mêmes commandes sont reconnues par `scripts/render.mjs`, qui produit la
% vue de lecture du site. Une macro ajoutée ici doit l'être aussi là-bas,
% faute de quoi le rendu échouera bruyamment — ce qui est voulu : un rendu
% silencieusement faux est pire qu'un rendu absent.

\ProvidesPackage{grothendieck}[2026/08/08 Transcriptions du fonds Grothendieck]

\usepackage{amsmath,amssymb,amsthm}
\usepackage{mathrsfs}
\usepackage{stmaryrd}
% Les diagrammes commutatifs sont partout dans ce fonds : c'est la langue
% même de la géométrie algébrique de Grothendieck, et une transcription qui
% les rendrait en prose ne transcrirait rien.
\usepackage{tikz-cd}
% Français par défaut — c'est la langue des feuillets — mais l'anglais est
% chargé aussi : la traduction et le résumé sont des documents anglais, et la
% typographie française (espaces avant les deux-points) y serait une faute.
% Ces documents basculent par \selectlanguage{english} en tête de corps.
\usepackage[english,french]{babel}
\usepackage{fontspec}
\usepackage{geometry}
\geometry{a4paper,margin=25mm,marginparwidth=28mm}
\usepackage{marginnote}
\usepackage{xcolor}
% ulem, not soul. `soul`'s \st and \ul are built on a character-by-character
% scan that XeLaTeX's font machinery breaks: the document fails with an
% undefined control sequence rather than degrading. ulem does the same two jobs
% and is Unicode-engine safe. `normalem` stops it hijacking \emph, which the
% transcriptions use for Grothendieck's own emphasis.
\usepackage[normalem]{ulem}
\usepackage{hyperref}
\hypersetup{colorlinks=true,linkcolor=black,urlcolor=black,pdfborder={0 0 0}}

% --- Métadonnées du lot -----------------------------------------------------
% Reprises telles quelles par le script de rendu, qui les affiche en tête de
% la vue de lecture. Elles fixent ce que le fichier prétend couvrir : sans
% elles, une transcription flotte sans point d'attache dans un fonds de seize
% mille feuillets.

\newcommand{\folder}[1]{\gdef\grothFolder{#1}}
\newcommand{\batch}[1]{\gdef\grothBatch{#1}}
\newcommand{\leaves}[2]{\gdef\grothFirst{#1}\gdef\grothLast{#2}}
\newcommand{\foldertitle}[1]{\gdef\grothFoldertitle{#1}}
\gdef\grothFolder{?}\gdef\grothBatch{?}\gdef\grothFirst{?}\gdef\grothLast{?}\gdef\grothFoldertitle{}

% --- Le résumé --------------------------------------------------------------
%
% Il ouvre la lecture modernisée, sous le titre « Résumé », et il est composé
% en petit. Ce n'est pas un ornement : c'est le seul passage du document qui
% ne soit pas des mathématiques, et un lecteur doit pouvoir le reconnaître
% avant de l'avoir lu — puis le sauter s'il connaît déjà le sujet, ou s'y
% arrêter s'il ne le connaît pas. Le filet qui le suit marque la reprise du
% corps.

\newenvironment{resume}
  {\small\setlength{\parskip}{0.45em}}
  {\par\medskip\hrule\medskip}

% --- Mots-clés ---------------------------------------------------------------
%
% En anglais, en clôture du résumé de la lecture modernisée : le vocabulaire
% moderne sous lequel chercher ce que les feuillets construisent. Le manifeste
% du site (`npm run manifest`) les extrait pour étiqueter la cote entière —
% cette ligne est la source unique des tags, il n'y a pas de fichier à côté.
\newcommand{\keywords}[1]{\par\smallskip\noindent
  {\footnotesize\textsc{Keywords} --- \textit{#1}}\par}

% --- L'appareil critique ----------------------------------------------------

% Le numéro de feuillet, en marge. C'est l'ancre qui permet de régler un
% désaccord sur une formule en regardant *un* feuillet plutôt qu'une chemise
% de six cents. Le site s'en sert aussi pour tourner les pages du fac-similé
% quand on fait défiler la transcription.
\newcommand{\leaf}[1]{%
  \marginnote{\footnotesize\color{gray!70}\textsf{#1}}%
  \label{leaf:#1}\ignorespaces}

% Illisible : on ne devine pas. Un mot inventé qui se lit comme les autres est
% le pire résultat possible.
\newcommand{\ill}{\textcolor{red!60!black}{[\,\ldots\,]}}

% Lecture douteuse : le mot est donné, et signalé comme incertain.
\newcommand{\uncertain}[1]{\uline{#1}}

% Ajout de l'éditeur — une lettre manquante, un mot sous-entendu.
\newcommand{\add}[1]{\textcolor{blue!50!black}{[#1]}}

% Biffé par Grothendieck lui-même. Ce qu'il a rayé fait partie du document :
% une rature dit souvent où la pensée a changé de direction.
\newcommand{\struck}[1]{\sout{#1}}

% Note du transcripteur, sur l'état matériel du feuillet.
\newcommand{\note}[1]{\par\smallskip{\small\color{gray!60!black}\textit{[#1]}}\par\smallskip}

% Note marginale de Grothendieck, distincte du corps du texte.
\newcommand{\marginal}[1]{\par\smallskip\noindent\hspace{1em}%
  {\small\color{gray!60!black}#1}\par\smallskip}

% --- Titre ------------------------------------------------------------------

\AtBeginDocument{%
  \begin{center}
    {\large\bfseries Fonds Alexandre Grothendieck --- cote n\textsuperscript{o}~\grothFolder}\\[2pt]
    {\itshape\grothFoldertitle}\\[4pt]
    {\small feuillets \grothFirst--\grothLast\ (lot \grothBatch)}\\[2pt]
    {\footnotesize\color{gray!60!black}Transcription établie d'après le fac-similé numérisé
     par l'Université de Montpellier.\\
     Les lectures incertaines sont soulignées, les illisibles notées [\,\ldots\,],
     les ajouts entre crochets.}
  \end{center}
  \vspace{1em}}

\endinput


\folder{45}
\batch{4}
\pages{61}{79}
\dating{1961-[vers 1970]}
\watermark{Édition de démonstration}
\foldertitle{Groupes semi-simples : notes manuscrites (s.d.), tapuscrit annoté (1961)}

\begin{document}

\page{61}

\note{la page porte deux fois le même schéma de flèches entre catégories de
formes. Le premier, en haut, est barré de grands traits courbes et de
hachures ; on en donne les éléments lisibles : un cadre « Spéc Bil typ
$\to$ Spéc.\ Bil unimod », puis « Bil typ.\ », « Bil unimod », « Bil »,
« Quad unimod », « Quad typ.\ » (ce dernier mot surchargé), « Quad », avec
$\text{Bil} \rightleftarrows \text{Quad}$, et « Mod » au bas, où aboutissent
trois flèches, dont une venant d'un « Alt » biffé ; « Spéc Bil typ /
Spéc Bil unimod », hachuré ; « Alt $\leftarrow$ Spéc.\ Alt » ; « N.B. ».
Le second, en bas, est propre ; on le redonne ci-dessous}

\[
(ax - by)^2 + (cx + dy)^2 = x^2 + y^2
\qquad
\left\lbrace
\begin{array}{l}
a^2 + c^2 = 1 \\
b^2 + d^2 = 1 \\
2(-ab + cd) = 0
\end{array}
\right.
\]
\note{en haut à droite, dans la partie barrée ; la fin du second membre
est surchargée, la lecture $x^2 + y^2$ est douteuse, de même que le signe
de la dernière ligne ; le second membre commençait par « $ax^2 + \ldots$ », biffé ; à côté, un mot souligné deux fois, non lu (\ill{})}

\begin{tikzcd}[column sep=small, row sep=small, nodes={font=\scriptsize}]
\text{Spéc Bil typ.} \arrow[r] \arrow[d] & \text{Bil typ} \arrow[d] & \text{Quad unim.} \arrow[dl] \arrow[d] & \text{Alt} \arrow[ddd] & \text{Spéc Alt} \arrow[l] \arrow[ddd] \\
\text{Spéc Bil unim.} & \text{Bil unim.} \arrow[d] & \text{Quad typ.} \arrow[d] & & \\
 & \text{Bil} \arrow[r, bend left=15] & \text{Quad} \arrow[l, bend left=15] & & \\
 & & \text{Mod} & \text{Mod unim.} \arrow[l] & \text{Spéc Mod unim.} \arrow[l]
\end{tikzcd}

\note{schéma redessiné. « Spéc Bil typ.\ / Spéc Bil unim.\ » et « Bil typ /
Bil unim.\ / Bil » sont chacun entre crochets. Plusieurs traits n'ont pas
de pointe visible ; on les a orientés vers le bas, comme ceux qui en ont
une. Les deux traits partant de « Alt » et de « Spéc Alt » descendent en
biais ; on les fait aboutir à « Mod unim.\ » et à « Spéc Mod unim.\ », au
début desquels ils s'arrêtent. Isolés au bas de la page : « M.\ » et « N »}

\page{63}

\[
e \to SO(2m) \to O(2m) \to \mathbb{Z}/2\mathbb{Z} \to e
\]
\note{sous $SO(2m)$, « centre $\mu_2$ », puis, sous une accolade, « centre
$\mu_2$ » (la première écriture surchargée) ; de $O(2m)$ part un trait
pointillé vers « produit semi-direct » ; sous $\mathbb{Z}/2\mathbb{Z}$, une
flèche verticale, et « dét $\to$ » vers un $\mathbb{G}_m$ noirci}
\[
e \to SO(2m+1) \to O(2m+1) \to \mu_2 \to e
\]
\note{même disposition : « centre $\mu_2$ » (la première fois surchargé),
« prod.\ semi-direct », et sous $\mu_2$ une flèche verticale vers
$\mathbb{G}_m$, marquée « dét »}

Pour tout entier $m$, on a

\begin{tikzcd}
e \arrow[r] & \mu_2 \arrow[r] & \mathrm{Pin}(m) \arrow[r] & O(m) \arrow[r] & e \\
e \arrow[r] & \mu_2 \arrow[r] \arrow[u, no head, "\Vert" description] & \mathrm{Spin}(m) \arrow[r] \arrow[u] & SO(m) \arrow[r] \arrow[u, hook] & e
\end{tikzcd}

\note{à côté de la flèche verticale du milieu, le mot « Spinoriel »,
biffé ; l'égalité de gauche est deux traits verticaux}

\struck{$x$, $x'$} \quad \ill{} car.~$2$ \qquad $\lambda$ \quad
\note{en haut à droite ; à gauche des formules suivantes, un petit carré
dessiné}
\begin{gather*}
(x, y) \longmapsto (\lambda x + u(y),\ y) \\
(\lambda(x) + u(y))^2 + Q_0(y) \\
\lambda^2 x^2 + u(y)^2 = x^2 \\
u(y)^2 = 0 \\
\lambda^2 = 1 \\
u(y)^2 = \\
\underline{\lambda^2 = 1} \qquad \mu_2 \\
(\mu_2\ \ \alpha_2\ \ \alpha_2\ \ \ldots\ \ \alpha_2)
\end{gather*}
\note{la première formule est surchargée ($\lambda$ et $u(y)$ écrits
par-dessus d'autres signes) ; le « $Q_0$ » est lu sur un signe peu formé.
La dernière ligne est entre grandes parenthèses}

\struck{$\lambda\ \ldots$}
\begin{gather*}
(\xi, y) \longmapsto (\lambda \xi + u(y),\ y) \\
\lambda^2 = 1 \qquad u(y)^2 = 0 \\
\lambda' \lambda \xi + \lambda' u(y) + u'(y) \\
(\lambda', u')(\lambda, u) = (\lambda' \lambda,\ \lambda' u + u') \\
\mu_2 .
\end{gather*}

\page{65}

Formes bilinéaires \uncertain{sym.}\ \uncertain{forment} un anneau

Les formes alternées sur modules sur cet anneau

Les formes quadratiques également
\note{les deux lignes sont réunies par une accolade ; « alternées » et « modules » sont d'une lecture douteuse}

F \ill{} \ill{}

$H^1(S, \mu_2)$ \note{le $H^1$ surcharge un autre signe ; une flèche en
part vers la ligne suivante}

\begin{quote}
Bil \quad $\lambda$-anneau \struck{\ill{}}

Alt \quad Module

Quad \quad Module
\end{quote}
\note{les deux dernières lignes sont réunies par un crochet marqué
« $\pi$ », d'où part une flèche vers la première ; à gauche, un signe
hachuré et un petit schéma biffé portant $\alpha$ et $\beta$}

Bil $+$ Alt \quad $\lambda$-module gradué par $\mathbb{Z}/2\mathbb{Z}$

muni d'un élément privilégié d'\uncertain{augm.}~$2$

muni d'un élément privilégié d'\uncertain{augm.}~$1$
\note{« Bil $+$ Alt » est séparé du reste par un trait vertical ; les deux
dernières lignes sont en face de « Alt Module » et « Quad Module » ;
« privilégié » est surchargé la première fois}

$A$ \quad $M$ \qquad $f(xy)z = x f(yz)$

\struck{$\lambda$-anneau} \qquad \struck{$Q(x \wedge y) =$} \qquad
$f \leadsto Q$ \qquad $(Q \otimes f') \, (\Sigma x_i$
\note{la dernière formule est inachevée et son $\Sigma$ surchargé ; au-dessus,
un signe isolé, peut-être $\mathfrak{g}^2$}

\[
f(x \wedge y, x \wedge y) =
\begin{vmatrix}
f(x, x) & f(x, y) \\
f(y, x) & f(y, y)
\end{vmatrix}
= Q(x, x) Q(y, y) - f(x, y)^2
\]
\[
Q(x \wedge y) = Q(x, x) Q(y, y) - f(x, y)^2
\]
\note{la matrice est écrite entre grandes parenthèses}

\struck{$\alpha(Q \otimes$} \qquad \struck{$\alpha(f \otimes f') =$} \qquad
$\alpha(f) = f \otimes \alpha(1)$ \qquad $\beta($

\note{le tiers droit de la page est écrit la feuille tournée d'un quart
de tour ; on le donne remis droit, de haut en bas}
\begin{gather*}
e_1 e_2 = e_1^2 = 0 \\
(x e_1 + y e_2)(x' e_1 + y' e_2) = y y' \\
(e_1 + e_2)(e_1 + e_2) \qquad e_1 + e_2 - \ldots + e_n \qquad
\Sigma x_i^2 \\
e_1\ \ \underbrace{e_2\ e_3 - \ldots}\ \ e_1 - e_2 \\
e_1,\ e_1 - e_2,\ e_1 - e_3,\ \ldots,\ e_1 - e_n \\
e'_i e'_j = 0 \quad i, j \neq 1 \\
a_{11} x_1 x'_1 + \sum_{1 \leq i < j} a_{ij} (x_i x'_j + x_j x'_i)
\qquad (e_1 - e_i)(e_1 - e_j) = 1 \\
a_{11} x_1 x'_1
\end{gather*}
\note{en tête, « $x\,x^2$ » biffé ; à côté, un petit dessin d'axes marqué « ad $t$ ». Les lignes
$(x e_1 + \ldots) = y y'$ et $(e_1 + e_2)(e_1 + e_2)$ paraissent barrées ; les indices de la somme sont
surchargés ($i \neq j$ corrigé en $1 \leq i < j$). Sous
$e_2 e_3 - \ldots$, une accolade avec deux mots non lus, et à gauche
« \ill{} » ; « $i, j \neq 1$ » est d'une lecture douteuse}

\page{67}

\[
0 \longrightarrow \Sigma^2(\check E) \longrightarrow \mathrm{Symm}^2(E)
\longrightarrow \Sigma^2(\check E) \longrightarrow \mathrm{Symm}^2(E)
\]
\note{au-dessus des deux premiers termes, les rangs
$\frac{n(n+1)}{2}$ ; sous chaque flèche du milieu, un triangle passant par
$\check E^{(2)}$, la descente marquée « surj », la montée « inj » ; au bout,
$\times \otimes x$ biffé}

\[
\textstyle\sum a_{ij} x_i x_j \leadsto \sum a_{ij} x_i \otimes x_j
\]
\[
0 \to \mathrm{Alt}^2(\check E) \to \Sigma^2(\check E) \to
\mathrm{Symm}^2(\check E) \to \mathrm{Alt}^2(\check E) \to 0
\]
\note{sous la flèche du milieu, un triangle passant par $\check E^{(2)}$}
\[
0 \to \check E^{(2)} \to \mathrm{Symm}^2(\check E) \to \Sigma^2(\check E)
\to \check E^{(2)} \to 0
\]
\note{sous la flèche du milieu, un triangle passant par
$\mathrm{Alt}^2(\check E)$}

\[
\boxed{\mathrm{Im} = \mathrm{Ker} = \mathrm{Alt}^2(\check E)
\qquad \Sigma^2(\check E) \rightleftarrows \mathrm{Symm}^2(\check E)
\qquad \mathrm{Im} = \mathrm{Noyau} = \check E^{(2)}}
\]
\note{dans un cadre ; les deux flèches entre $\Sigma^2$ et
$\mathrm{Symm}^2$ sont deux arcs}

$\mathbb{R}$ \quad $\mathbb{C} = \mathbb{R}[i]$ \quad
$K = \mathbb{R}[i, j, k] = \mathbb{C} \mathbin{\bar\otimes} \mathbb{C}$
\quad \struck{$M_2(\mathbb{C})$} \quad \struck{$M_2(K)$}
\note{colonne en haut à droite, séparée du reste par un trait ; le signe
au-dessus du $\otimes$ est peu formé}

\ill{} $+$ \uncertain{Alternée} \ill{} $+$ \uncertain{bilin.}\ \ill{}.
\qquad $\Sigma X_i$ \qquad
$f \otimes f'(x \otimes x', y \otimes y') = f(x, y)\, f'(x', y')$
\note{le premier $\otimes$ de la formule surcharge un autre signe}

\struck{Alt} \quad Forme quadratique

Forme bilinéaire symétr.

Forme \struck{symplectique} alternée
\note{les deux premières lignes sont réunies par une accolade, la troisième
en a une à elle seule ; « alternée » est écrit au-dessus de
« symplectique » biffé}

\begin{quote}
Forme \struck{symplectique} alternée $\otimes$ Forme \struck{symplectique}
alternée $\to$ forme \struck{symplectique} \struck{\uncertain{bilinéaire}}
bilin.\ \uncertain{sym.}

forme \struck{\uncertain{bilinéaire}} $\otimes$ forme bil.\ $\to$ forme
altern.

forme bil $\otimes$ forme bil $\to$ forme bil.

Forme quadratique $\otimes$ forme bil.\ $\to$ forme quadr.

forme quadratique $\otimes$ forme alternée $\to$ forme alternée

Forme quadratique $\otimes$ forme quadratique ?!..?
\end{quote}
\note{les trois premières lignes sont réunies par une accolade à gauche ;
les corrections de la première sont écrites au-dessus des mots biffés, et
la dernière est d'une lecture douteuse ; le mot biffé de la deuxième
ligne n'est pas sûr}

\page{69}

$E^{2n}(S) = {}$ \uncertain{ens.}\ des classes de couples $(E, F)$, $E$ et
$F$ \uncertain{formes} quadratiques admissibles de rangs
\struck{\ill{}} \uncertain{pairs} tels que
$\operatorname{rang} E - \operatorname{rang} F = {}$
\note{« de rangs » est surchargé et « pairs » est écrit au-dessus d'un
mot biffé ; la ligne s'arrête sur le signe $=$ ; l'exposant $2n$ est
d'une lecture douteuse}

\[
x_1^2 + \cdots + x_n^2 + \struck{x_1 x_2}\ y_1 y_2 + y_3 y_4 + \cdots +
y_{2m-1} y_{2m}
\]

$E^{n,m} = {}$ \uncertain{ens.}\ des classes \uncertain{d'isom.}\ de formes
quadratiques \uncertain{régulières} de type $(n, m)$

\[
\coprod_{n, m \geq 0} E^{n,m}(S) = \bar E(S)
\]
\uncertain{est} muni d'une structure de \emph{\uncertain{monoïde}}
\emph{bigradué}

\struck{Soit} [N.B. Si $2 \to$ \ill{}, alors
$E^{n,m} \simeq E^{n+2m, 0}$

si $2 = 0$, \struck{A \ill{}} alors $E^{n,m} \simeq E^{1,m}$

\struck{\ill{}} \ill{} \ill{}, \ill{} \ill{} \ill{} \ill{} \ill{}
\struck{\ill{}}

\ill{} \ill{} $E^{\alpha, 0}$ ou $E^{1, \beta}$, ou $E^{0, \gamma}$
\ldots\ ]
\note{« $2 \to$ \ill{} » est peut-être « 2 inversible » ; l'exposant
$n+2m$ est d'une lecture douteuse ; sous le second $E^{1,m}$, un petit
signe}

$E(S)$ \uncertain{groupe} \uncertain{associé} au \uncertain{monoïde}
\[
E(S) \xrightarrow{\ \varepsilon_1 \times \varepsilon_2\ } \mathbb{Z}
\times \mathbb{Z}
\]
\note{à gauche, une formule barrée d'un grand trait oblique, « $\coprod E^{n,m}(S) / n + \ldots$ » ; au-dessus,
un signe $=$ vertical la rattache à $E(S)$}

$E^{n,}$
\note{inachevé}

\subsection*{Commentaires sur l'article de A.~Delzant sur les formes quadratiques}

\page{71}

\marginal{22.7.961}
\note{date de sa main, en haut à droite, à l'encre ; le deuxième chiffre
surcharge un autre. Pages 71 à 74 : texte dactylographié, avec ses
corrections à l'encre. Les lettres grecques ($\varphi$, $\xi$, $\lambda$,
$\mu$) et le $\wp$ de la page 74 sont ajoutés à l'encre dans les blancs
laissés par la machine}

Du point de vue rédaction;

a) Il est très confusant de faire semblant de traiter en même temps la
\uncertain{carc.}\ $\neq 2$ et la car.~$2$. Pour l'instant, même les
structures qui entrent en jeu sont d'espèce différente, et le traitement
des deux cas est vraiment différent. Si donc tu n'arrives pas à une
théorie commune satisfaisante, laisse tomber la car.~$2$, ou mets dans un
numéro à part ce que tu sais en dire.

b) Contrairement à ce que suggère la lecture de l'article, le rôle des
$\lambda$-structures dans la théorie est pratiquement nul, dû au fait sans
doute que la filtration naturelle de l'anneau des classes de formes
quadratiques n'est autre que celle définie par les puissances de l'idéal
d'augmentation. On récupère trivialement la $\lambda$-structure sur la
forme explicite dudit anneau, et il suffit de le faire à la fin.

Du point de vue math., je suggère comme programme minimum raisonnable le
suivant, qui est à portée immédiate:

1) \emph{Structure de l'anneau} $Q(A)$ \emph{des classes de formes
quadratiques sur} $A$, ou pour l'instant $A$ va être un corps de
caractéristique $\neq 2$. Pour $a \in A^{*}$, soit $\varphi_a$ la classe de
la forme quadratique $a\xi^2$, elle ne dépend que de $a \bmod A^{*2}$,
d'où un homomorphisme de
\[
D(A) = A^{*}/A^{*2} = H^1(A, \mathbb{Z}/2\mathbb{Z})
\]
dans $Q(A)^{*}$, notée aussi $a \to \varphi_a$, d'où un homomorphisme
(*) \[
\varphi : \mathbb{Z}(D) \longrightarrow Q(A)
\]
où $\mathbb{Z}(D)$ est l'algèbre du groupe $D = D(A)$, à coéfficients dans
$\mathbb{Z}$. On prouvera alors le théorème suivant:
\note{« $Q(A)^{*}$ » : un exposant tapé est surfrappé ; « (*) » et
« $\varphi :$ » sont ajoutés à l'encre}

\emph{L'homomorphisme précédent est surjectif}, \emph{son noyau est le
sous-groupe engendré par les} $e_a + e_b - e_{a'} - e_{b'}$, \emph{pour}
$a, b, a', b' \in D$ \emph{tels que} $\varphi_a + \varphi_b = \varphi_{a'}
+ \varphi_{b'}$. D'après un théorème de Witt, cette dernière relation
équivaut aussi à dire que les formes des deux membres ont mêmes
invariants, ce qui s'écrit plus simplement
\[
a + b = a' + b' \quad , \quad a.b = a'.b'
\]
où $D$ est écrit additivement, et où le produit $a.b$ est le cup-produit,
à valeurs dans
\[
B = H^2(A, \mathbb{Z}/2\mathbb{Z})
\]
\marginal{?}
\note{un trait vertical au crayon dans la marge gauche, avec un point
d'interrogation, en face de « relation équivaut aussi à dire » jusqu'à
« à valeurs dans »}

\page{72}

considéré comme sous-groupe du groupe de Brauer formé des éléments annulés
par 2. Donc \emph{la struuture de} $Q(A)$ \emph{est connue quand on connait
les groupes} $H^i(A, \mathbb{Z}/2\mathbb{Z})$ \emph{pour} $i = 1, 2$,
\emph{soient} $D$ \emph{et} $B$, \emph{et le cup-produit} $D \times D \to
B$. L'énoncé précédent montre \struck{en effet} même que le noyau de
l'homomorphisme surjectif (*) est aussi le sous-groupe engendré par les
éléments
\[
e_{a+u+v} - e_{a+u} - e_{a+v} + e_a
\]
pour les éléments $a, u, v$, dans $D$ tels que $uv = 0$.
\note{« même » est écrit à l'encre au-dessus de « en effet » biffé}

On en conclut en particulier que le gradué associé à $Q(A)$ est isomorphe
au quotient de l'algèbre symétrique de $D$ sur $\mathbb{Z}$, par l'idéal
engendré par les éléments d'ordre 2 de la forme $ab - a'b'$, où $a, b, a',
b'$, sont comme ci-dessus, ou ce qui revient au, même, engendré par les
éléments $uv$ de $\mathrm{Symm}^2(D)$ pour $u, v$ $D$ tels que $u.v = 0$
dans $B$ (N.B. il faudrait éclaircir si $\mathrm{grad}^2 Q(A) \to B$ est
\uncertain{tjs} injectif, en d'autres termes si le noyau de
$\mathrm{Symm}^2(D)$ $B$ est engendré par des éléments de la forme $uv$).
\note{le mot devant « injectif » est tapé puis surchargé, peut-être à
l'encre ; lecture douteuse}

On conclut du théorème de structure que si $H^{\cdot}$ est un anneau gradué
à degrés positifs, alors la donnée d'un homomorphisme \struck{additif}
d'anneaux augmentés de $Q(A)$ dans $\mathbb{Z} \times (1 + \hat
H^{+})$ équivaut à la donnée d'un homomorphisme
\[
W^1 : D \longrightarrow H^1
\]
tel que $W^1(a) W^1(b) = W^1(a') W^1(b')$ lorsque $a, b, a', b'$ sont comme
ci-dessus ou ce qui revient au même, tel que $W^1(u) W^1(v) = 0$ si
$u.v = 0$ dans $B$; d'ailleurs $\mathrm{grad}\, Q(A)$ joue un rôle
universel pour ces propriétés, i.e.\ la donnée de $W$ équivaut aussi à la
donnée d'un homomorphisme d'anneaux gradués
\[
\mathrm{grad}\, Q(A) \to H^{\cdot} .
\]
\note{« d'anneaux augmentés » est écrit à l'encre au-dessus de
« additif » biffé ; « $\mathbb{Z} \times ($ » et « $)$ » sont ajoutés à
l'encre autour de « $1 + \hat H^{+}$ »}

Le choix le plus naturel pour $H^{\cdot}$, du point de vue "topologique",
semble
\[
H^{\cdot}_{1} = H^{\cdot}(A, \mathbb{Z}/2\mathbb{Z}) ,
\]
muni du cup-produit, ce qui nous donne des classes de
\emph{Stiefel-Whitney}
\[
W^i : Q(A) \to H^i(A, \mathbb{Z}/2\mathbb{Z}) .
\]
Serre suggère de considérer aussi une définition plus fine des classes de
Stiefel-Whitney, en remplaçant l'extension séparable algébrique maximale de
$A$ par la sous-extension réunion des sous-extensions finies de rang une
puissance de 2; cela donnerait déjà un $H^2$ "plus fin", i.e.\ s'envoyant
dans le $H^2$ habituel \ldots
\note{l'indice de $H^{\cdot}_{1}$ est surfrappé}

\page{73}

Bien entendu, $Q(A)$ est un $\lambda$-anneau, et sa $\lambda$-structure se
lit sur la structure d'anneau, puisqu'on aura
\[
\lambda_t(\varphi_a) = 1 + \varphi_a t ,
\]
D'autre part, un homomorphisme d'anneaux augmentés $Q(A) \to \mathbb{Z}
\times (1 + \hat A^{+})$ est automatiquement un $\lambda$-homomorphisme. Je
ne crois pas qu'il y ait lieu de dire plus des $\lambda$-structures (sauf
peut-être en remarque que la filtration mise sur $Q(A)$ est celle déduite
de sa $\lambda$-structure).

2) \emph{Cas où} $A$ \emph{est un anneau local}, de corps résiduel $k$ tel
que car $k \neq 2$. Il semble bien que tout ce qui précède marche dans ce
cas; bien entendu, $H^{\cdot}(A)$ désigne la cohomologie du "groupe
fondamental" de $\mathrm{Spec}(A)$, au sens de mon séminaire. Il faudrait
alors tout rédiger dans le cadre: formes quadratiques sur un anneau local.
N.B. une telle forme est dite non dégénérée si la réduite est non
dégénérée; en principe, il devrait être possible de façon assez triviale
de relever les constructions habituelles sur $k$, en des constructions
sur $A$.

Bien entendu, il serait très agréable de dépasser le programme minimum et
d'arriver à une théorie des classes de Stiefel-Whitney pour des Modules
quadratiques sur un (pré)schéma quelconque $S$, en supposant au besoin que
les caractéristiques résiduelles sont $\neq 2$. A titre de test, on peut
toujours essayer de prouver l'invariance de la définition triviale
relative à un Module quadratique qui splitte en modules quadratiques de
rang 1.
\note{« (pré) » devant « schéma », et « quadratique », « quadratiques »
après « Module » et « modules », sont ajoutés à l'encre ; un trait d'encre
passe sur « qui splitte » et sur « de rang », sans qu'on puisse dire s'il
les biffe}

Comme prime au lecteur, on peut en tous cas donner la compatibilité de la
définition purement algébrique des classes de Stiefel-Whitney, et de la
définition topologique, lorsque $A$ est l'anneau local d'un point $s$ d'un
schéma de type fini $S$ sur $\mathbb{C}$. Un module quadratique sur $A$
définit un module quadratique sur un voisinage de $S$, donc $S$ si on veut,
i.e.\ un fibré prinnipal homogène algébrique sous le schéma en groupes
orthogonal. Par Gaga il lui correspond un fibré principal analytique, donc
topologique, d'où des classes de Stiefel-Whitney topologiques dans
$H^{\cdot}(S_{\mathrm{top}}, \mathbb{Z}/2\mathbb{Z})$. Or on a un
homomorphisme canonique, que tu définiras:
\[
H^{\cdot}(S_{\mathrm{top}}, \mathbb{Z}/2\mathbb{Z}) \to H^{\cdot}(A,
\mathbb{Z}/2\mathbb{Z}) ,
\]
et on voit aussitôt que les images des classes caractéristiques
topologiques sont les classes caractéristiques algébriques, (car il suffit
de le vérifier en rang 1).
\note{« d'un point $s$ » est écrit à l'encre au-dessus de « local d'un »,
et « $S$ » après « de type fini »}

\page{74}

3) \emph{Essayer de trouver des définitions raisonnables en car.}\,2,
\emph{ou de préférence indépendantes de la caractéristique}. Une théorie
vraiment satisfaisante devrait marcher sur tout préschéma $S$, et
réciproquement. Je suis loin d'avoir des idées claires sur la question, et
me borne à des remarques vagues:

a) Lorsqu'on est sur un corps de caractéristique \struck{2, on est}
quelconque, sans doute obligé de tenir compte des formes quadratiques
dégénérées, et de définir un anneau $Q'(k)$ comme défini par les classes
de formes quadratiques, dégénérées ou non. [Dans le cas d'un corps de
caractéristique $\neq 2$, ce serait $Q(k) + \mathbb{Z}$, où la
multiplication est donnée par $(f, m).(g, n) = (fg,\ \deg(f) n + \deg(g) m
+ mn)$, le facteur en rabiot $\mathbb{Z}$ correspondant aux formes nulles,
sur des vectoriels de rang quelconque].
\note{« quelconque » est écrit à l'encre au-dessus de « 2, on est »
biffé ; les crochets sont à l'encre}

b) En caractéristique quelconque, garder à l'esprit la distinction entre
le shhéma en groupes $(\mathbb{Z}/2\mathbb{Z})_S$ et le shhéma en groupes
$(\mu_2)_S$ sur $S$, le dernier étant le noyau de la puissance deuxième
dans le groupe multiplicatif $\mathbb{G}_{m\,S}$. Si $S$ est local
d'anneau $A$, alors la suite exacte de cohomologie donne $H^1(S, \mu_2) =
A^{*}/A^{*2}$, $H^2(S, \mu_2) = {}_2H^2(A, \mathbb{G}_m)$, éléments annulés
par 2 dans le groupe de Brauer, par contre $H^1(S, \mathbb{Z}/2\mathbb{Z}) =
A/\wp A$ lorsque $2A = 0$. Il faut donc s'attendre à trouver des classes de
Stiefel-Whitney de type mixte, faisant à la fois intervenir les deux
groupes voisins de coéfficients. Cette remarque se précise quand on
regarde les schémas en groupes orthogonaux, sur $\mathbb{Z}$ disons: en
dimension paire, c'est l'extensions de $\mathbb{Z}/2\mathbb{Z}$ par un
schéma en groupes à fibres connexes (le schéma spécial-orthogonal), par
exemple en dim 2 c'est le produit semi-direct de $\mathbb{Z}/2\mathbb{Z}$
par $\mathbb{G}_m$ sur lequel il opère par symétrie; en dimension impaire,
c'est une extension de $\mu_2$ par le groupe spécial-orthogonal, par
exemple en dimension 1, c'est $\mu_2$. D'où suivant la parité, comme
premier invariant "discret" d'un fibré principal orthogonal, une 1-classe
de groupe $\mathbb{Z}/2\mathbb{Z}$, ou $\mu_2$.
\note{le $\wp$ de $A/\wp A$ est ajouté à l'encre}

\subsection*{Formes quadratiques (Exposés de Kneser à l'IHES)}
\note{titre de sa main en tête de la page 75}

\page{75}

\marginal{(1) \emph{Généralités}}
\note{le numéro 1 est entouré}

$A$ anneau commutatif

$E(A)$ anneau des formes quadratiques "non dégénérées"

\quad [i.e.\ à discriminant non diviseur de $0$] sur modules projectifs de
type fini

$F(A)$ anneau des formes quadratiques "unimodulaires"

\quad [i.e.\ : discriminant inversible]

Si $A$ est un corps, $E(A) = F(A)$

\emph{Invariants pour} $F(A)$
\begin{itemize}
  \item[1°)] Rang \quad $\varepsilon : F(A) \longrightarrow \mathbb{Z}$
  \quad hom d'anneaux
  \item[2°)] Discriminant \quad $d : F(A) \longrightarrow A^{*}/A^{*2}$
  \quad hom de groupes
  \item[3°)] Classe de Clifford \quad $c : F(A) \longrightarrow
  {}_2\mathrm{Br}(A)$ \quad application
\end{itemize}

La définition de $c$ par Kneser se fait via la définition d'abord pour
\uncertain{les} modules quadratiques [en distinguant le cas pair, où on
prend $Cl(E)$, et impair, où on prend $Cl^{+}(E)$], puis on note que si
$E$ et $E'$ ont \struck{un rang} discriminant \uncertain{unités}
\uncertain{dans} \ill{} $c(E + E') = c(E) + c(E')$, donc on peut définir
$c$ au moins sur les sous-groupe $F^2(A)$ noyau de $\varepsilon, d$
\ldots
\note{les mots de liaison de ce paragraphe sont d'une écriture rapide ;
« via », « d'abord », « distinguant », « impair » sont des lectures
probables, non sûres}

[N.B. Lorsque $A$ est local de car.\ rés.\ $\neq 2$, on peut donner une
définition des $w_i$ de Stiefel-Whitney, redonnant $d$ et $c$ \ldots

\page{76}

$\varepsilon, d, c$ donnent une filtration de $F(A)$
\[
\left.
\begin{array}{l}
F(A)/F_1(A) \simeq \mathbb{Z} \\
F_1(A)/F_2(A) \simeq A^{*}/A^{*2} \\
F_2(A)/F_3(A) \subset {}_2\mathrm{Br}(A)
\end{array}
\right\rbrace
\]
dans le cas $A$ un corps de car.\ $\neq 2$
\note{une flèche va de ${}_2\mathrm{Br}(A)$ à « sous-groupe des éléments
\uncertain{décomposables} »}

\struck{$F_1$} \quad A-t-on en tout cas pour tout corps de car.\
\uncertain{arbitraire}
\[
\left\lbrace
\begin{array}{l}
F_1(A)/F_2(A) \simeq H^1(A, \mathbb{Z}/2\mathbb{Z}) \\
F_2(A)/F_3(A) \simeq H^2(A, \mathbb{Z}/2\mathbb{Z})
\end{array}
\right.
\]
\note{le reste du feuillet ne porte que la transparence de la page 75}

\page{77}

\marginal{(2)}
\note{numéro entouré}

\emph{Exemples élémentaires} (des corps) \quad (tous ces corps de car
$\neq 2$)
\note{« des corps » est écrit dans une bulle au-dessus de la ligne}

\begin{itemize}
  \item[a)] \emph{corps alg.\ clos} $k$ \quad $F(k) \simeq \mathbb{Z}$ \quad
  [On a $F_1(k) = 0$]
  \item[b)] \emph{Corps fini} $k$ \quad $F(k) \simeq \mathbb{Z} +
  \mathbb{Z}/2\mathbb{Z}$ \quad [On a $F_2(k) = 0$]
  \item[c)] \emph{Corps} $p$-\emph{adique} $k$ \quad $F(k) \simeq
  \mathbb{Z}$ \ill{} \uncertain{un} \uncertain{p.\ gr.}\ fini
  \uncertain{annulé} par 4 \quad [On a $F_3(k) = 0$]
  \item[d)] \emph{Corps} $\mathbb{R}$ \quad $F_1(k) \simeq \mathbb{Z}$
  \uncertain{qui} \ill{} \quad $F_2(k) = 2\mathbb{Z}$ \quad $F_3(k) =
  4\mathbb{Z}$ \quad \struck{$F$}
  \item[e)] \emph{Corps valué complet} à corps résiduel alg.\ clos
  \struck{\ill{}} \quad $F_1(k) \simeq \mathbb{Z} + \underbrace{k^{*}/k^{*2}}_{\mathbb{Z}/2\mathbb{Z}}$
  si car.\ résiduelle est $\neq 2$, \uncertain{grâce} \ill{} \ldots
  \quad [On a $F_2(k) = 0$]
\end{itemize}
\note{« b) » surcharge « a) » ; l'indice de $F_1$ en e) est peu formé}

\marginal{(3)}
\note{numéro entouré ; devant le titre, un signe noirci}

\emph{Corps de nombres}

\emph{Th.\ de Hasse} \quad $F(K) \longrightarrow \sum_{p} F(K_p)$ \quad
\emph{injectif}
\note{« injectif » est souligné deux fois}
$F_1(K)/F_2(K) \longrightarrow \sum_p K_p^{*}/K_p^{*2}$ \quad
injectif [\uncertain{grâce} à \uncertain{Hasse} \ill{}]

$F_2(K)/F_3(K) \longrightarrow \sum_p \mathbb{Z}/2\mathbb{Z}$ \quad
injectif [car $\mathrm{Br}(K) \subset \sum_p \mathrm{Br}(K_p)$]

$F_3(K) \subset \sum_{p} F_3(K_p) \subset \mathbb{Z}^r$ \quad ($p$ infini
réel) \quad $r$ nb de places réelles
\note{une flèche relie « \uncertain{sous-groupe} des \uncertain{familles}
d'éléments de somme nulle » à $F_2(K)/F_3(K)$ ; « \uncertain{nulle} » est
douteux}

Le sous-groupe de torsion de $F(K)$ est le noyau de
$F_1(K) \xrightarrow[\text{aux places réelles}]{\text{indices}}
\mathbb{Z}^r$, il est annulé par 4.
\note{les mots au-dessus et au-dessous de la flèche sont d'une lecture
douteuse}

\page{78}

\begin{tikzcd}[column sep=small, row sep=small]
0 \arrow[r] & E'(A) \arrow[r] \arrow[d, "2"] & E(A) \arrow[r] \arrow[d, no head] & E(K) \arrow[r] \arrow[d, no head] & 0 \\
0 \arrow[r] & \sum E'(\hat A_p) \arrow[r] & \sum E(\hat A_p) \arrow[r] & \sum E(K_p) \arrow[r] & 0
\end{tikzcd}

\note{les deux dernières verticales sont de simples traits ; le « 2 » de
la première est un exposant, lecture douteuse}

\[
\mathrm{Pic}(A)/2\,\mathrm{Pic}(A) \subset H^2(A, \mu_2)
\]
\note{écrit au-dessus de $F'(A)$, auquel le relie un signe $\simeq$
vertical ; au-dessus de $F(K)$, « $E(K)$ » avec un signe $=$}

\begin{tikzcd}[column sep=small, row sep=small]
0 \arrow[r] & F'(A) \arrow[r] & F(A) \arrow[r] \arrow[d] & F(K) \arrow[d] & \\
 & 0 \arrow[r] & \sum F(\hat A_p) \arrow[r, "\text{inj}"'] & \sum F(K_p) &
\end{tikzcd}

\note{la première ligne se termine par un trait noirci}

$A$ anneau \ill{} \ill{} \ill{} \ldots
\note{en haut à droite, un petit dessin barré portant « $x^2 = a$ »,
« $x^2 - a + y^2 = 0$ » et « $x/t$ »}

\[
H^1(K, \mathbb{Z}/2\mathbb{Z}) \to \sum_p H^1(K_p, \mathbb{Z}/2\mathbb{Z})
\qquad
H^1(K, \mathbb{G}_m) \to
\]
\note{à gauche, « \struck{$\mathbb{Z}$}$/2\mathbb{Z}$ » ; tout ce qui suit,
jusqu'au trait horizontal, est pris dans de grandes courbes qui le
barrent}

\struck{$0 \to T \to J_T \to J/T \to 0$}

\struck{$K$ : \ $0 \to L^{*} \to J_L \to C_L \to 0$}

\struck{$0 \to L^{*} \otimes X \to J \otimes X \to C \otimes X \to 0$}

\struck{$\sum_p H^1(G_p, L_p^{*} \otimes X) \to$}

\struck{$H^1$ \quad $H^2(G, \ldots) \to H^2($}

\struck{$H^2(G, K^{*}) \to H^2(\ ,J)$}
\note{lecture très incertaine des indices ; les quatre suites sont dans
un grand cadre courbe, à droite}

\uncertain{anneau} \uncertain{complet} de \uncertain{val.}\
\uncertain{discrète}

$E'(V)$ \uncertain{est} \ill{} : \ill{} \ill{} \ill{} d'une \ill{} \ill{}
de $E(k)$, avec $k$ le corps résiduel
\note{« anneau complet de val.\ discrète » est au-dessus, relié par une
flèche courbe à $V$}

\end{document}
